% ID: SEL-D0026-OB-E02
% ESTADO: validada. S027, cotejada con examen y criterios.
\textbf{a)} El dominio es $\mathbb R\setminus\{0\}$.
En $0$ la función no está definida y sus límites laterales
son infinitos de signos opuestos. En $2$ los dos límites y
el valor son $2$. Es continua en todo su dominio.
Las derivadas en los tramos abiertos son $-4/x^2$ y $2x-2$.
En el empalme dan pendientes laterales $-1$ y $2$.
Así, es derivable en $\mathbb R\setminus\{0,2\}$.

\textbf{b)} Para $x>0$, el primer tramo tiene derivada negativa
y el segundo positiva. Decrece en $(0,2]$, crece en $[2,\infty)$
y tiene un mínimo en $x=2$. Las segundas derivadas son $8/x^3>0$
en $(0,2)$ y $2>0$ en $(2,\infty)$. Ambos tramos son convexos
(cóncavos hacia arriba). Al pasar por $2$ la pendiente salta
hacia arriba, así que el empalme conserva la convexidad global
en $(0,\infty)$.
